\clearpage
%\chapter*{\centerline{\LARGE{General Introduction}}}
\chapter*{{\emph\LARGE{General Conclusion and Perspectives}}}
%\{General Introduction}
\markboth{General Conclusion and Perspectives}{}
\addcontentsline{toc}{chapter}{General Conclusion and Perspectives}
\pagestyle{fancy}
\fancyhf{}
\fancyhead[HL]{General Conclusion and Perspectives}
\cfoot{\thepage}



%%%%%%%%%%%%%%%%%%%%%%%%%%%%%%%%%%%%%%%%%%%
\lettrine[lines=2,loversize=0.1]{\textbf{\textit{T}}}his graduation project was carried out within Be Wireless Solutions as 
part of the attainment of the National Engineering Diploma in Electronics 
and Embedded Systems. 
 work presented in this report addresses a concrete engineering challenge 
faced by Be Wireless Solutions: the absence of a common, reusable firmware 
base for industrial IoT deployments that previously required complete 
redevelopment for each new hardware configuration. The IOBOX platform was 
designed to eliminate that problem by providing a single, modular embedded 
system capable of adapting to heterogeneous sensor and communication 
environments without structural changes to the firmware.
intended to serve as a shared and reusable base across future BWS deployments.


Throughout this report, we presented the different phases of our work. We 
began by establishing the general context of the project, introducing the 
host company, its activities, and the industrial need that motivated this work. 
We then defined the software architecture of the system, and explained the design choices that guided our 
implementation.

The development phase covered the construction of a layered firmware 
structure, from low-level peripheral drivers up to the application logic. 
We implemented sensor acquisition, industrial communication protocols, 
signal processing, and data persistence, all organized within a clean and 
extensible architecture. The platform was then subjected to a series of 
hardware validation tests that confirmed the correct behavior of every 
major component and interface on the target microcontroller.

This project required several months of sustained work, combining embedded software development, protocol integration, and 
systematic testing. It gave us the opportunity to face real engineering 
challenges  from low-level driver debugging to architectural decisions 
and to develop a deliverable that holds genuine industrial value for the company.

Despite the results achieved in this first version, we believe that several 
improvements and extensions could strengthen the platform further, and we 
set them as perspectives:

\begin{itemize}
    \item Migrate from the cooperative scheduler to a lightweight real time 
    operating system to better handle time critical tasks.
    \item Extend the communication stack with support for additional 
    industrial protocols such as Modbus TCP .
    \item Add remote firmware update capability allowing system parameters 
    and feature flags to be adjusted at runtime rather than at compile time.
    \item Extend the platform with a secure communication layer, adding 
    authentication and data encryption to meet industrial cybersecurity requirements.
\end{itemize}
